% !TeX program = xelatex
\documentclass[11pt,a4paper]{article}
\usepackage[margin=23mm]{geometry}
\usepackage{fontspec}
\setmainfont{Latin Modern Roman}
\usepackage{amsmath,amssymb,mathtools}
\usepackage{xcolor,enumitem,booktabs,longtable,array,xurl,hyperref}
\definecolor{muted}{HTML}{53636C}
\hypersetup{colorlinks=true,urlcolor=blue,linkcolor=blue}
\setlength{\parindent}{0pt}
\setlength{\parskip}{5pt}
\newcommand{\R}{\mathbb R}
\newcommand{\E}{\mathbb E}
\newcommand{\N}{\mathbb N}
\newcommand{\ind}{\mathbf 1}
\newcommand{\Var}{\operatorname{Var}}
\newcommand{\Cov}{\operatorname{Cov}}
\newcommand{\supp}{\operatorname{supp}}
\DeclareMathOperator*{\argmin}{arg\,min}
\DeclareMathOperator*{\argmax}{arg\,max}
\newcommand{\entrymeta}[3]{{\sffamily\small\color{muted}\textbf{\texttt{#1}}\quad|\quad #2\par #3\par}}
\newcommand{\Problemref}[1]{\texttt{#1}}
\newcommand{\JELref}[2]{\texttt{#2} (JEL \texttt{#1})}
\begin{document}
\section*{Coordinating multiple tax instruments}
\textbf{Problem ID:} \texttt{H21-58}\quad \textbf{JEL:} \texttt{H21}\par
% BEGIN ECONOMICS STATEMENT
\label{problem:OP-0471}
{\sffamily\small\color{muted}Original subject group: Taxation and social insurance\par}
\label{definitions:problem:30007753}
\textbf{Audit classification:} Published research agenda. The integrating-instruments agenda is verified. Added a finite-horizon equilibrium definition, common numeraire, instrument equivalence criterion, and supremum/attainment distinction.
\subsection*{Definitions and precise research target}
Fix a finite horizon \(H\), a population of private household types \(\theta\sim F\), observable characteristics \(x(\theta)\), production and trading technologies, initial assets, an information structure and terminal conditions. A policy \(p=(T_y,T_k,t_c,T_f)\in\mathcal P\) specifies the complete labor- and capital-income schedules, commodity taxes and family transfers at each date. The admissible set \(\mathcal P\), administrative resource cost \(C(p)\), tax bases, remittance dates and any behavioral misperceptions are declared in advance. Every schedule uses only legally observable variables. A competitive equilibrium \(e\in\mathcal E(p)\) comprises household plans, firm plans, prices, receipts and market-clearing conditions at every date, with sequential budget feasibility and the specified terminal conditions. Exclude policies without an equilibrium and assume bounded integrable utilities.

Let \(U_\theta(e)\) be the declared lifetime welfare index and \(G_\theta\) the fixed social valuation. For every admissible pair \((p,e)\), require \(\theta\mapsto G_\theta(U_\theta(e))\) to be measurable and absolutely integrable with respect to \(F\), and put \(W(e)=\int G_\theta(U_\theta(e))\,dF(\theta)\). Bounded integrability of \(U_\theta\) alone does not supply this requirement. If outcomes are uncertain, the definition of \(U_\theta\) specifies their welfare evaluation. Let \(B(p,e)\) be discounted receipts net of transfers and administration. Declare dates and, where relevant, states, their probability law and information, valuation factors or state prices, government borrowing and interim budget constraints, terminal public obligations, and the disposition of receipts exceeding the funding requirement. These financing and resource rules form part of equilibrium feasibility. Discount all public flows using the specified valuation factors and date-zero numeraire; the numeraire alone does not determine the factors. Declare whether the comparison \(B(p,e)\geq R\) concerns a scalar evaluation of contingent receipts (specifying that evaluation), or a statewise comparison (specifying every-state or almost-sure feasibility and the corresponding scalar or contingent requirement \(R\)). Fix this interpretation across policies and instrument classes. For funding requirement \(R\), define
\[
 V_{\mathcal P}=\sup_{p\in\mathcal P}\inf_{e\in\mathcal E(p)}\int G_\theta(U_\theta(e))\,dF(\theta),
 \qquad B(p,e)\geq R\quad\forall e\in\mathcal E(p).
\]
For any declared instrument class \(\mathcal Q\subseteq\mathcal P\), let \(\mathcal F_{\mathcal Q}=\{p\in\mathcal Q:\mathcal E(p)\ne\varnothing,\ B(p,e)\geq R\ \forall e\in\mathcal E(p)\}\), using the declared revenue interpretation. The supremum defining \(V_{\mathcal Q}\) ranges over \(\mathcal F_{\mathcal Q}\). Report an empty feasible class separately. For nonempty classes the value may be extended real; no uniform welfare bound across policies and equilibria is assumed. Write \(w(p)=\inf_{e\in\mathcal E(p)}W(e)\). Fix a baseline policy \(p_b\in\mathcal F_{\mathcal P}\) before comparing reforms. Under the worst-equilibrium criterion a feasible reform \(p\) has a positive welfare gain precisely when \(w(p)>w(p_b)\); this compares worst-equilibrium levels and does not assert improvement for every pair of equilibria. If a selected-equilibrium comparison is used instead, declare one selection \(s(p)\in\mathcal E(p)\) on all policies with nonempty equilibrium sets, define \(w^s(p)=W(s(p))\) and \(V^s_{\mathcal Q}=\sup_{p\in\mathcal F_{\mathcal Q}}w^s(p)\), and use \(w^s(p)>w^s(p_b)\) for a positive gain. Revenue feasibility remains robust over all equilibria in either welfare comparison. Selected values are distinct from the displayed robust values. Characterize joint reforms satisfying the applicable positive-gain and revenue conditions; characterize when a declared smaller instrument class \(\mathcal P_0\subseteq\mathcal P\) has value \(V_{\mathcal P_0}=V_{\mathcal P}\), or separately \(V^s_{\mathcal P_0}=V^s_{\mathcal P}\) under a declared selection. The numerical welfare loss is \(V_{\mathcal P}-V_{\mathcal P_0}\) when both values are finite, with the analogous convention for selected values. Otherwise report the individual values and feasibility status without forming an undefined difference; finite gains are likewise reported only when the compared levels are finite. Let \(a(e)\) denote the complete equilibrium allocation and define \(A(p)=\{(a(e),B(p,e)):e\in\mathcal E(p)\}\). For correspondence-based comparisons, coincidence of allocations and receipts means \(A(p)=A(p')\); for selected comparisons it means \((a(s(p)),B(p,s(p)))=(a(s(p')),B(p',s(p')))\). Declare which comparison is intended. Policies are allocation-equivalent only if their allocations and net public receipts coincide in this declared sense; matching one labor wedge does not establish equivalence. Neither equality of instrument-class values nor preservation of utility alone establishes allocation equivalence. Existence of an optimizer requires additional compactness and continuity arguments. This is a newly scoped policy problem. It is a parameterized characterization agenda: a particular theorem or numerical instance must state its model instance or quantified model class and its success criterion, without treating any benchmark as a solution of the general agenda.
\subsection*{Variants and assumption changes}
\begin{enumerate}
\item \textbf{Static separability (editorial benchmark).} Use one period, identical homothetic consumption preferences weakly separable from labor, perfect observability of the declared tax bases and unrestricted nonlinear earnings taxes. As benchmark inputs, list all observable and legally taxable variables, explicitly stating whether ability and effort remain private, and specify the full and smaller instrument classes and which instruments are removed or restricted. Perfect observability of declared bases does not itself make every private characteristic observable or taxable. Test instrument redundancy under these explicit restrictions and the declared welfare and revenue criterion, rather than assuming redundancy in a heterogeneous dynamic economy.
\item \textbf{Restricted instruments (editorial).} Limit earnings taxes to a declared finite number \(K\geq1\) of brackets and add commodity taxes. Specify whether bracket thresholds are fixed inputs or policy choices, together with threshold domains, continuity requirements and admissible rates and intercepts. Declare the commodity tax bases, price convention and admissible rate domain, and define uniform taxation by a common rate under that convention. Use the same earnings-tax class in the full class and the smaller class \(\mathcal P_0\), which restricts commodity taxation to the declared uniform form. Compute \(V_{\mathcal P}-V_{\mathcal P_0}\) when both values are finite and both feasible classes are nonempty; otherwise report values and feasibility separately. A numerical instance specifies \(\varepsilon>0\) and a certified interval of width at most \(\varepsilon\) containing the finite loss, with any reported candidate policies feasible under the declared criterion. This differs from unrestricted income taxation and does not presume a positive loss.
\item \textbf{Behavioral welfare (editorial).} Specify distinct decision and welfare utilities and an observation model for misperceived taxes. Correction of behavioral distortions is a possible policy motive alongside redistribution; these specifications alone imply neither optimizer existence nor that permitted instruments achieve correction.
\end{enumerate}
\subsection*{References and source audit}
Kaplow's 2022 manuscript section 7 provides the integration framework; section 9 discusses extensions. Journal metadata are verified on the 2024 AEA page. The exact multi-instrument minimax question is editorial. In the cited Harvard manuscript, section 7 begins on printed p. 69, section 7.3 discusses commodity-tax benchmarks on printed pp. 76--77, and section 9 is on printed pp. 89--90; the latter page range does not locate all sections listed in reference 1. The association of the working-paper DOI with this exact July 2022 Harvard version remains unverified. The journal metadata check does not verify the journal full text.
\begin{enumerate}\raggedright
\item\hypertarget{definitions:source:30007753:1}{} Louis Kaplow. 2022. \emph{Optimal Income Taxation}. Olin Discussion Paper 1083, July 2022; sections 3.2, 7, 8.2 and 9, printed pp. 89–90.
\url{https://laweconcenter.law.harvard.edu/wp-content/uploads/2024/11/Kaplow_1083.pdf}.
DOI: \url{https://doi.org/10.3386/w30199}.
\item\hypertarget{definitions:source:30007753:2}{} Louis Kaplow. 2024. \emph{Optimal Income Taxation}. Journal of Economic Literature 62(2), 637–738; publisher abstract.
\url{https://www.aeaweb.org/articles?id=10.1257/jel.20221647}.
DOI: \url{https://doi.org/10.1257/jel.20221647}.
\end{enumerate}

\subsection*{Common conventions: Source mathematical conventions}

These conventions apply to each entry unless explicitly replaced.

\textbf{Models and identification.} A primitive model \(m\) specifies all probability laws, feasible choices, information, equilibrium and measurement rules used by its target. The admissible class \(\mathcal M\) is fixed independently of outcomes. If \(P_O(m)\) is its observable probability law and \(\tau(m)\) the target, its identified set at observable law \(P\) is
\[
 \mathcal I_\tau(P)=\{\tau(m):m\in\mathcal M,\ P_O(m)=P\}.
\]
Point identification means that this set is a singleton. An empty set signals incompatibility of data and assumptions. Coordinate bounds need not describe the joint set. Bounds are sharp only if every retained target is attainable or arbitrarily approachable by compatible models. Finite bounds require explicit restrictions on unobserved quantities; unrestricted confounding, hidden wealth, extrapolation or arbitrary equilibrium selection can yield uninformative sets.

\textbf{Causal interventions.} A potential outcome \(Y_i(p)\) is the outcome under a complete policy regime \(p\), including timing, financing, information and market spillovers. The target population and units are fixed before treatment. With interference, use the complete assignment vector or a justified exposure mapping; individual no-interference notation alone cannot identify an economy-wide intervention. A difference of mean potential outcomes does not require the cross-world joint distribution. Conditioning on post-treatment migration, survival, enrollment or employment changes the estimand and needs separate justification.

\textbf{Statistical statements.} An estimator, confidence set, forecast or test specifies a sampling law, model class and loss/coverage criterion. Uniform \((1-\alpha)\)-coverage means \(\inf_{m\in\mathcal M}\Pr_m\{\tau(m)\in C_n\}\geq1-\alpha\), or an explicitly asymptotic version. The whole parameter space trivially has coverage, so informativeness requires a stated width, risk or power objective. A forecasting improvement is relative to a named benchmark, information set and evaluation population.

\textbf{Equilibrium and optimization.} An equilibrium comprises feasible plans, optimality, beliefs where needed, budget constraints and all market/resource clearing. The primitives or a stated benchmark define each correspondence \(\mathcal E(p;m)\). Nonemptiness is assumed or proved, never inferred just from compact policy sets. A selected-equilibrium target uses a declared selection rule; a robust target quantifies over all permitted equilibria. Use suprema or infima unless attainment is established. An \(\varepsilon\)-optimal policy is feasible and within \(\varepsilon\) of the finite optimum value. Report an empty feasible set separately.

\textbf{Welfare and accounting.} Normative utility, interpersonal normalization, social weights, numeraire, discounting and the use of public receipts are primitives. Behavioral utility need not coincide with the welfare criterion. Household welfare includes ownership income and rents once; adding profits again to compensating variation generally double-counts them. A monetary equivalent is defined by an explicit transfer and price convention, with monotonicity and range conditions. Average effects, mechanism contributions, transfers and welfare losses are different targets. Infinite sums require convergence and dynamic feasible plans require terminal or no-Ponzi conditions.

\textbf{Source versions.} A working-paper date, revision date and journal publication year are distinguished. A DOI for a journal version is not automatically the DOI of an earlier manuscript. Locators refer to the stated version and printed pages where specified. A metadata-only check does not verify a claim about a full-text conclusion. Every entry's source subsection narrows the provenance claim accordingly.
% END ECONOMICS STATEMENT
\end{document}
